\subsection{Zielerreichung und Projektreflexion}\label{sec:evaluation_reflexion}

Zwischen zwei Festlegungen dieses Berichts besteht eine Spannung. Die Einleitung begründet die gesellschaftliche Relevanz über \ac{SDG} 12.3, also eine Verringerung der Lebensmittelverschwendung auf breiter Ebene. Die Zielgruppe ist dagegen bewusst verengt, weil die Gamification-Mechanik vor allem bei bereits motivierten Personen wirkt. Beide Aussagen passen zusammen, wenn Resteria zunächst als Angebot für diese engere Nische wirkt und die breitere \ac{SDG}-12.3-Wirkung erst über Skalierung und die in~\autoref{sec:kernergebnisse_und_ausblick} genannten Kooperationen entsteht.

Die Kombination aus Wettbewerbsanalyse und literaturbasierter Herleitung der Gamification-Funktionen war für diesen Projektrahmen angemessen, weil sie Markt und verhaltenswissenschaftliche Fundierung ohne eigene Nutzerdaten abdeckt. Zwei Stellen erzwangen eine Kursänderung. Der zunächst naheliegende Community-Zuschnitt mit Forum und Nutzergruppen wich zeitlich begrenzten Challenges (siehe~\autoref{sec:community_und_feedback}), und eine Namenskollision zwang zur Korrektur des Arbeitstitels.

Der Weg zu diesem Konzept folgte einer erkennbaren Abfolge: Zuerst legte diese Arbeit mit der Zero-Waste-Nische den Fokus fest, dann folgten die Sichtung der sechs Plattformen im Juli 2026 und die Literaturrecherche zu Nudges und Gamification, aus denen sich Funktions- und Gamification-Konzept ableiteten und anschließend in die UI/UX-Mockups übersetzt wurden.

Aus der zweiten dieser Kursänderungen lässt sich eine Lehre ziehen: Der Arbeitstitel \enquote{Restlos} kollidierte mit der bestehenden Zero-Waste-Initiative \enquote{Restlos Glücklich} (Berlin, seit 2015, identisches Themenfeld). Ein früher Stresstest deckte das auf, sodass die Korrektur zu \enquote{Resteria} nur eine Umbenennung kostete – ähnlich günstig wie die in~\autoref{sec:iteratives_design} beschriebenen frühen Überarbeitungsrunden an den Mockups.
